% 3.7 — Deux systèmes, un décideur
\section{Deux systèmes, un décideur}
\label{sec:deux-os}

Les agents ne travaillent pas de la même façon sous Linux et sous Windows : sous Windows,
Claude Code dispose d'un outil \cmd{PowerShell} à côté de son outil \cmd{Bash} (Git Bash), et
Codex écrit du PowerShell. Les chemins s'écrivent \cmd{C:\textbackslash x} ou \cmd{/c/x}, il
n'y a pas de ptrace, et les appels système n'ont rien de commun. La conception sépare donc
ce qui \emph{doit} être commun de ce qui \emph{ne peut pas} l'être
(figure~\ref{fig:os}).

\begin{figure}[htbp]
\centering
% Figure 5 — Linux et Windows : un seul décideur, deux filets d'exécution.
\begin{tikzpicture}[
    col/.style={boite,text width=57mm,font=\sffamily\footnotesize},
    titre/.style={font=\sffamily\bfseries\small,text=encre},
  ]
  % Partie commune
  \node[accent,text width=72mm,minimum height=10mm] (eng) at (0,0)
     {Moteur partagé : simulation, effets, impact\sous{même table de règles, même \texttt{policy.py}}};
  \node[boite,draw=encre,line width=1.1pt,text width=72mm,below=4mm of eng] (dec)
     {\textbf{\texttt{decide.py}} --- un seul décideur\sous{mêmes garanties sur les deux systèmes}};
  \draw[fleche] (eng) -- (dec);

  % Frontaux
  \node[boite,text width=40mm] (fp) at (-4.75,1.9) {Frontal POSIX\sous{\texttt{model.py} ; Git Bash : \texttt{/c/x} $\rightarrow$ \texttt{C:\textbackslash x}}};
  \node[boite,text width=40mm,dashed] (fw) at (4.75,1.9) {Frontal PowerShell (C3)\sous{lieur de paramètres ; \statutCthree}};
  \draw[fleche] (fp.south) |- (eng.west);
  \draw[fleche] (fw.south) |- (eng.east);

  % Colonnes OS
  \node[titre] (tl) at (-4.75,-2.9) {Linux};
  \node[titre] (tw) at (4.75,-2.9) {Windows};
  \node[col,below=2mm of tl] (l1) {Agents : Claude Code (Bash), Codex,\\piège bash (\texttt{BASH\_ENV}, \texttt{trap DEBUG})};
  \node[col,below=2mm of l1,fill=wepwnclair,draw=wepwnfonce] (l2) {\textbf{Garde d'appels système}\sous{\texttt{ptrace} : arrêt à l'entrée de \texttt{unlink}, \texttt{rename},\\\texttt{open(O\_TRUNC)}\ldots{} ; copie \emph{puis} laisse faire}};
  \node[col,below=2mm of l2] (l3) {Activation : réécriture par le \textit{hook}\\\texttt{demo\_cli sh -g <id> <cmd>} (sans \texttt{allow})};
  \node[col,below=2mm of l3] (l4) {Plan de contrôle : appel annulé\\(\texttt{orig\_rax = -1}) + reçu ESCALATE};

  \node[col,below=2mm of tw] (w1) {Agents : Claude Code (Bash \& PowerShell),\\Codex (PowerShell)};
  \node[col,below=2mm of w1,fill=wepwnclair,draw=wepwnfonce] (w2) {\textbf{Montage WinFsp}\sous{système de fichiers en espace utilisateur : \texttt{rename}\\(remplacement), \texttt{cleanup}(DELETE), \texttt{set\_file\_size}}};
  \node[col,below=2mm of w2] (w3) {Activation : le projet \emph{est} le montage\\(\texttt{C:\textbackslash p} $\rightarrow$ données dans \texttt{C:\textbackslash p.real} verrouillé)};
  \node[col,below=2mm of w3] (w4) {Remontage à l'ouverture de session :\\tâche élevée durcie (W-SEC)};

  \draw[fleche] (dec.south) -- ++(0,-0.5) -| (l1.north);
  \draw[fleche] (dec.south) -- ++(0,-0.5) -| (w1.north);

  \node[zone,fit=(tl)(l1)(l4)] (zl) {};
  \node[zone,fit=(tw)(w1)(w4)] (zw) {};
  \node[etiquette,text width=120mm,anchor=north] at ([yshift=-2mm]0,0 |- zw.south) {Les deux filets couvrent le \emph{projet} seulement, et la raison de la décision le dit.};
\end{tikzpicture}

\caption{Linux et Windows : un seul décideur, deux filets d'exécution}
\label{fig:os}
\end{figure}

\begin{description}[style=nextline,leftmargin=0.6cm,font=\sffamily\bfseries]
  \item[Commun]
    Le moteur, la table de règles, la politique, le décideur, le format des reçus, la
    récupération et \cmd{undo}. Une ligne POSIX tapée dans Git Bash sous Windows passe par le
    même moteur que sous Linux ; ses chemins \cmd{/c/\ldots} sont traduits en
    \cmd{C:\textbackslash\ldots} (partie M2.1), et un chemin POSIX absolu dont la racine MSYS
    n'est pas connue (\cmd{/tmp}, \cmd{/usr}) est marqué non compris plutôt que deviné.
  \item[Propre à chaque système]
    Le filet d'exécution et son activation, la protection du plan de contrôle pendant
    l'exécution, et le frontal de lecture : POSIX d'un côté, PowerShell de l'autre (C3,
    section~\ref{sec:c3}).
\end{description}

\subsection{Linux : pourquoi ptrace}
Le filet Linux est un \emph{garde d'appels système} fondé sur ptrace
(\fichier{syscall\_guard.py}). Il lance la commande comme processus suivi et s'arrête à
l'\emph{entrée} de chaque appel système ; pour un appel destructeur (\cmd{unlink},
\cmd{unlinkat}, \cmd{rmdir}, \cmd{rename*} sur une destination existante, \cmd{truncate},
\cmd{open} avec \cmd{O\_TRUNC} ou en écriture) sur un fichier préexistant du projet, il copie
le fichier \emph{puis} laisse l'appel se faire. Pour un fichier du plan de contrôle, il annule
l'appel.

Trois mécanismes étaient envisageables :
\begin{itemize}
  \item \textbf{FUSE} : monter un système de fichiers sous le projet. Il a été écarté parce que
    le garde d'appels système fait déjà la même chose --- voir chaque destruction dans le
    projet avant qu'elle n'ait lieu. De plus, un montage FUSE suppose de placer le projet
    derrière un point de montage servi par un démon permanent, alors que ptrace ne s'attache
    qu'à la commande concernée.
  \item \textbf{seccomp avec notification} : plus rapide, car seuls les appels filtrés
    réveillent le superviseur. Le code le désigne comme la cible de production du filet ; le
    prototype actuel utilise ptrace, qui donne la main à l'entrée de chaque appel et permet de
    lire les registres et la mémoire du processus suivi. \acompleter[confirmer la raison du choix de ptrace pour le prototype]
  \item \textbf{ptrace} : retenu. Ses limites sont assumées et rappelées au
    chapitre~\ref{ch:securite} : architecture x86-64, coût d'un arrêt par appel système, et
    pas de vision au-delà du projet.
\end{itemize}

\paragraph{Activation.} Le filet doit s'exécuter \emph{à la place} de la commande. Le
\textit{hook} Claude Code réécrit donc la commande en \cmd{demo\_cli sh -g <id> <commande>}
via \cmd{updatedInput} --- \emph{sans} répondre \cmd{allow}, pour que l'invite de permission
de l'utilisateur continue de s'afficher (partie B1b) --- et dit à la personne et au modèle ce
qui s'exécute réellement. Un \cmd{demo\_cli sh} tapé par l'agent lui-même est déballé et jugé
sur ce qu'il cache ; un processus déjà suivi par un traceur étranger refuse de s'exécuter
sans surveillance.

\subsection{Windows : pourquoi un montage WinFsp}
Windows n'a pas de ptrace. L'endroit où toutes les routes vers un fichier convergent --- API
Win32, \cmd{ntdll}, appels système directs --- est le système de fichiers. \demo{} y place un
système de fichiers en espace utilisateur WinFsp (\fichier{fsmount.py}) : le projet
\cmd{C:\textbackslash projet} devient le point de montage, et les données vivent dans
\cmd{C:\textbackslash projet.real}, verrouillé par une ACL réservée aux Administrateurs et à
SYSTEM. Le montage s'exécute avec élévation, l'agent sans : l'agent ne peut pas contourner le
montage en écrivant directement dans le répertoire réel (\fichier{protect.py}). Le filet
n'exige donc aucune réécriture de commande : tout ce qui touche le projet passe par lui.

La logique de jugement (\fichier{fsguard.py}) est séparée de la plomberie WinFsp et testable
sous Linux. Ses règles viennent de journaux d'opérations \emph{observés} sous Windows~10,
qui ont contredit l'intuition sur trois points sur quatre : le remplacement atomique
(\cmd{rename} avec remplacement) détruit la destination sans aucun signal de suppression ---
c'est ainsi que l'outil \cmd{Write} de Claude Code enregistre chaque fichier ; la méthode
\cmd{overwrite()} n'est jamais appelée, les vraies réécritures passant par une réduction de
taille ; et seul \cmd{cleanup} avec le bit de suppression signale toutes les suppressions.

Un montage meurt avec la session. Une tâche planifiée à l'ouverture de session le relance ;
elle s'exécute avec élévation, ce qui en fait une cible : sa conception durcie (W-SEC) est
décrite à la section~\ref{sec:wsec}.

\subsection{Points de contrôle : utiles, mais optionnels}
Quand ni capture ciblée ni filet ne sont possibles, \demo{} peut copier tout le projet avant
la commande (\fichier{checkpoint.py}) --- « je ne sais pas ce que vous allez détruire, donc je
préserve tout ce que vous pourriez détruire ». Cette option est \textbf{désactivée par
défaut} : sur un dépôt avec un gros \fichier{.git}, la copie est lente, et « un garde-fou
qui devient lent sans qu'on le lui demande est un garde-fou qu'on désinstalle ». Elle est
bornée (256\,Mo et \num{25000} fichiers par défaut), refuse le répertoire personnel et la
racine, et ne remplace jamais une capture ciblée ni ne couvre une surface distante ou un
fichier hors du projet.
